% ==============================================
% ELEMENTO: CONCLUSÕES
% ==============================================
\section{Conclusões}
\label{sec:conclusoes}

Neste trabalho dimensionou-se um buffer de referência de 1,1~V para um DAC capacitivo a 250~MHz, formado por um seguidor
de saída $M_{11}$, uma réplica $M_1$ e um amplificador de dois andares com compensação indireta através da cascode. O
circuito final cumpre todas as especificações nominais: $A_0 = 63{,}1$~dB, $GBW = 107{,}9$~MHz, $PM = 66{,}5^\circ$, uma
queda máxima de 48,3~mV e um erro de \textit{settling} de 0,015\%, com 1,00~mA no amplificador e uma figura de mérito de
$1{,}54\times10^3$~V$^{-1}$.

A principal conclusão do trabalho diz respeito aos diferentes níveis de modelo usados. As expressões aproximadas,
obtidas a partir de um modelo de três nós, foram essenciais para o dimensionamento, porque mostram de forma direta que
parâmetro de cada transístor controla cada especificação. No entanto, cada simplificação afasta o modelo do circuito
real. As contas feitas com a tabela de referência previram um ganho de 49,3~dB, contra os 66,6~dB simulados, porque as
condutâncias de saída reais eram até 66\% inferiores às da tabela, obtida com um $V_{DS}$ diferente do do circuito. A
função de transferência exata, calculada em Python com os valores do ponto de funcionamento simulado, corrigiu o ganho e o
GBW, mas não a margem de fase, que tanto no dimensionamento inicial como no final ficou cerca de $15^\circ$ acima da
simulada. A razão é que os polos não dominantes não estão bem separados entre si nem do zero e formam um par complexo, pelo
que a hipótese $|p_1| \ll |p_2| \ll |p_3|$, base das expressões aproximadas, deixa de ser válida. Conclui-se que o modelo
aproximado serve para dimensionar e o modelo exato para compreender as diferenças, mas só a simulação permite validar a
estabilidade.

As iterações mostraram também quais os parâmetros que realmente limitavam o projeto. A margem de fase não era limitada
pela corrente do segundo andar, mas pela razão $\alpha = C_C/C_3$: aumentar $M_4$ e $M_5$ gastava corrente sem melhorar a
margem de fase, enquanto a correção de $C_P$ para $10\,C_{gs11}$ e o ajuste de $C_C$ a resolveram. Verificou-se ainda que
os graus de liberdade sem custo de corrente, como os multiplicadores dos transístores com a corrente fixa, devem ser
explorados antes de aumentar o consumo, já que alteram o $g_m/I_D$ e as capacidades sem custo para o FOM. Seguindo esta
ordem, o consumo do amplificador passou de cerca de 4~mA no dimensionamento inicial para 1~mA no circuito final.

Alguns compromissos só se tornaram visíveis no circuito completo. Reduzir a réplica ($K = 10$) poupa área e corrente, mas
baixa o polo do nó $V_{outR}$ e custa cerca de $7^\circ$ de margem de fase. O condensador $C_P$ protege $V_{ctr}$ da queda
provocada pelo $C_{DAC}$, mas faz parte de $C_3$ e afeta a estabilidade. No domínio do tempo, a análise mostrou que o erro
remanescente no fim do ciclo é de nível e não de velocidade: o erro total face a 1,1~V (0,071\%) é dominado pelo
\textit{offset} sistemático do amplificador, que não aparece na análise de pequeno sinal e pode ser corrigido com um
pequeno ajuste da corrente da réplica. A própria forma de medir (o instante de medida e a referência do erro) mostrou-se
determinante para a interpretação dos resultados.

Os restantes resultados confirmam a análise teórica: a impedância de saída (84,3~$\Omega$ em DC, com o polo a 178~MHz)
segue o modelo do seguidor de fonte, o ruído de saída (88,6~$\mu$V$_{rms}$) é dominado pelo ruído branco do amplificador
e não pelo seguidor, e o PSRR$^+$, de 65,3~dB em DC, desce para 9,3~dB à frequência de relógio, onde a malha já não tem
ganho, o que exigiria uma alimentação bem desacoplada.

Por fim, algumas margens do circuito final são reduzidas, nomeadamente 1,7~mV na queda inicial e cerca de 35~mV na
saturação de $M_3$ e $M_5$. O circuito cumpre as especificações no caso nominal, mas uma análise de \textit{corners} e de
Monte Carlo seria necessária para garantir o funcionamento com variações de processo, tensão e temperatura.